\selectlanguage{english}

In the General Relativity part, Lorentzian metrics are always taken with mostly-positive signature $ (1 , n-1) $, i.e. $ (-,+,\dots,+) $, therefore a vector $ X_p \in T_p\man $ is:
\begin{itemize}
  \item timelike if $ X_p^2 \equiv g_p(X_p , X_p) < 0 $;
  \item null if $ X_p^2 \equiv g_p(X_p , X_p) = 0 $;
  \item spacelike if $ X_p^2 \equiv g_p(X_p , X_p) > 0 $;
\end{itemize}


Greek indices are Lorentzian indices, while Latin indices are Euclidean indices (or generic $ \N_0 $- or $ \Z $-indices), and the Einstein summation convention on repeated indices is used unless otherwise specified.

Moreover, natural units are used, with $ c = \hbar = k_\text{B} \equiv 1 $, so that mass, energy, temperature, inverse length and inverse time all have the same units.
